% 研制总结报告插图：一次交互的数据流（主线 + 参数旁路 + AI/联机旁路）
\documentclass[tikz,border=8pt]{standalone}
\usepackage{xeCJK}
\IfFontExistsTF{Noto Sans SC}
  {\setCJKmainfont{Noto Sans SC}}
  {\IfFontExistsTF{Source Han Sans SC}
    {\setCJKmainfont{Source Han Sans SC}}
    {\setCJKmainfont{Microsoft YaHei}}}
\usetikzlibrary{positioning,arrows.meta,calc}

\definecolor{ink}{HTML}{1F2937}
\definecolor{cblue}{HTML}{2563EB}
\definecolor{cgreen}{HTML}{059669}
\definecolor{corange}{HTML}{D97706}
\definecolor{cred}{HTML}{B91C1C}
\definecolor{cgray}{HTML}{6B7280}

\begin{document}
\begin{tikzpicture}[x=1cm,y=1cm,font=\small,
    mod/.style={draw=#1!55!black, fill=white, rounded corners=3pt, align=center,
      inner sep=3pt, text=ink, font=\small, line width=0.6pt, minimum height=1.35cm},
    arr/.style={-{Latex[length=2mm]}, draw=ink!70, line width=0.8pt},
    darr/.style={arr, dashed, draw=cblue!70!black},
    lab/.style={font=\footnotesize, text=ink!85, inner sep=1pt},
    dl/.style={font=\footnotesize, text=cblue!70!black, inner sep=1pt}
  ]

  % 主线
  \node[mod=corange, minimum width=2.40cm] (src) at (1.50,0)
    {\normalsize\bfseries 倾斜 / 按键\\[1pt]\footnotesize IMU · KEY2};
  \node[mod=cblue, minimum width=2.90cm] (input) at (4.65,0)
    {\normalsize\bfseries 输入状态机\\[1pt]\footnotesize 去抖 · 定档 · 回差};
  \node[mod=cblue, minimum width=2.70cm] (game) at (8.25,0)
    {\normalsize\bfseries 规则核心\\[1pt]\footnotesize 预览 / 提交 / 结算};
  \node[mod=cgreen, minimum width=2.90cm] (ui) at (12.10,0)
    {\normalsize\bfseries 界面（LVGL）\\[1pt]\footnotesize 预览棋盘 · 结算};

  % 旁路
  \node[mod=cblue, minimum width=2.70cm] (cfg) at (8.25,2.35)
    {\normalsize\bfseries 参数配置\\[1pt]\footnotesize 10 项参数 · 三套配置};
  \node[mod=cblue, minimum width=2.50cm] (ai) at (8.25,-2.60)
    {\normalsize\bfseries AI 搜索\\[1pt]\footnotesize 快照进 · 动作出};
  \node[mod=cblue, minimum width=2.90cm] (net) at (12.10,-2.60)
    {\normalsize\bfseries 联机对战\\[1pt]\footnotesize 房主提交后广播};
  \node[mod=corange, minimum width=2.70cm] (peer) at (15.55,-2.60)
    {\normalsize\bfseries 对手板 / 旁观板\\[1pt]\footnotesize BLE 中央};

  % 主线箭头
  \draw[arr] (src.east) -- (input.west);
  \node[lab] at (2.95,0.88) {角度 · 按键};
  \draw[arr] (input.east) -- (game.west);
  \node[lab] at (6.50,0.88) {选择 / 确认};
  \draw[arr] (game.east) -- (ui.west);
  \node[lab] at (10.12,0.88) {棋盘快照};

  % 参数旁路（虚线：新开一局才生效）
  \draw[darr] (cfg.south) -- (8.25,0.68);
  \node[dl, anchor=east] at (8.12,1.28) {规则 · 搜索深度};
  \draw[darr] (cfg.west) -- (4.65,2.35) -- node[dl, left=2pt, pos=0.55]{缓冲参数} (4.65,0.68);

  % AI 旁路
  \draw[arr] (7.55,-0.68) -- node[lab, left, pos=0.5]{快照} (7.55,-1.93);
  \draw[arr] (8.45,-1.93) -- node[lab, right, pos=0.5]{动作} (8.45,-0.68);

  % 联机旁路
  \draw[arr] (9.42,-0.68) -- (9.42,-1.50) -- (12.10,-1.50)
    -- (12.10,-1.93);
  \node[lab] at (10.70,-1.30) {同步};
  \draw[arr] (net.east) -- (peer.west);

  % 图例
  \draw[arr] (0.40,-2.95) -- (1.35,-2.95);
  \node[font=\footnotesize, text=ink!85, anchor=west] at (1.45,-2.95) {一次交互的主线数据流};
  \draw[darr] (0.40,-3.45) -- (1.35,-3.45);
  \node[dl, anchor=west] at (1.45,-3.45) {参数：新开一局时写入};

\end{tikzpicture}
\end{document}
